\documentclass[11pt,a4paper]{article}

% ---------------------------------------------------------------------------
%  Fill these in before submitting (they appear on the title page and in §1).
% ---------------------------------------------------------------------------
\newcommand{\AuthorName}{G L Lohith}
\newcommand{\LiveURL}{[Vercel URL: add after deploy]}
\newcommand{\Passcode}{[APP\_PASSCODE: share in the submission form]}
\newcommand{\RepoURL}{[GitHub repository URL]}
\newcommand{\VideoURL}{[Demo video / GIF link]}
\newcommand{\SlidesURL}{\texttt{Turnout-slides.pdf} in the repository}
% ---------------------------------------------------------------------------

\usepackage[margin=2.2cm]{geometry}
\usepackage[T1]{fontenc}
\usepackage{libertinus}
\usepackage[scaled=0.86]{inconsolata}
\usepackage{microtype}
\usepackage{amsmath}
\usepackage{graphicx}
\usepackage{booktabs}
\usepackage{tabularx}
\usepackage{longtable}
\usepackage{array}
\usepackage{xcolor}
\usepackage{enumitem}
\usepackage{caption}
\usepackage{subcaption}
\usepackage{float}
\usepackage{titlesec}
\usepackage{fancyhdr}
\usepackage[numbers,sort&compress]{natbib}
\usepackage[hidelinks]{hyperref}
\usepackage{xurl}

\definecolor{accent}{HTML}{15803D}
\definecolor{muted}{HTML}{555555}
\hypersetup{colorlinks=true,linkcolor=accent,citecolor=accent,urlcolor=accent}
\graphicspath{{figures/}{../docs/img/}}
\setlist{itemsep=1pt,topsep=3pt,leftmargin=*}
\renewcommand{\arraystretch}{1.18}
\setlength{\parskip}{3pt}
\setlength{\emergencystretch}{3em}
\captionsetup{font=small,labelfont=bf}
\titleformat{\section}{\large\bfseries}{\thesection}{0.6em}{}
\titleformat{\subsection}{\normalsize\bfseries}{\thesubsection}{0.6em}{}
\pagestyle{fancy}\fancyhf{}
\lhead{\small\textsc{Turnout}: attendee growth command centre}
\rhead{\small RISA Growth APM, Part 2, Challenge A}
\cfoot{\small\thepage}

\newcommand{\rs}{Rs.\,}
\newcommand{\code}[1]{\texttt{#1}}
\newcommand{\lbl}[1]{\textsf{\small #1}}
\newcolumntype{L}{>{\raggedright\arraybackslash}X}

% Diagram: included when the image exists, otherwise a labelled placeholder.
\newcommand{\diagram}[3]{%
  \begin{figure}[H]\centering
  \IfFileExists{figures/#1}{\includegraphics[width=0.96\linewidth]{figures/#1}}{%
    \fbox{\parbox[c][4.6cm][c]{0.9\linewidth}{\centering\small\color{muted}
      Diagram to be supplied: \texttt{paper/figures/#1}\\[2pt]
      (generation prompt in \texttt{docs/diagram-prompts.md})}}}
  \caption{#2}\label{#3}\end{figure}}
% Screenshot from docs/img.
\newcommand{\shot}[3]{%
  \begin{figure}[H]\centering
  \includegraphics[width=0.94\linewidth]{#1}
  \caption{#2}\label{#3}\end{figure}}

\title{\vspace{-1.2cm}\textbf{Turnout}\\[4pt]
\large A grounded growth command centre for college-fest attendee growth\\[2pt]
\normalsize RISA Growth APM assignment, Part 2, Challenge A (Attendee Growth), Bengaluru}
\author{\AuthorName}
\date{November 2026}

\begin{document}
\maketitle
\thispagestyle{fancy}

\begin{abstract}
\noindent \textbf{Turnout} is a working growth system for \emph{CampusFest 2025}, the two-day festival in the assignment's dataset, which we set in Bengaluru, the district with the most colleges in India (1,148)~\cite{aishe2024}. It turns the supplied event-sourced workbook into a single analytics engine that \textbf{measures} the funnel against goal, \textbf{diagnoses} the weakest stage, \textbf{decides} which campaigns to scale or stop and where to move budget, and helps the team \textbf{manage} the resulting work. A tool-calling AI copilot answers questions only from that engine. It does arithmetic in a calculator, checks every number in its answer against tool output, never writes data without a confirmation, and shows a trace of exactly how each answer was produced. A voice assistant (Gemini Live) uses the same tools and can move people around the app.
On the original data, as of the 31 March decision point, the engine projects 2,887 registrations against a goal of 3,200 and 1,837 visitors against 2,200. It finds the bottleneck at identified$\to$registered (70.4\% against a plan-implied 88.9\%) and recommends moving \rs16,800 from two campaigns past the stop line to four cheaper ones (+157 registrations, +100 visitors). The copilot passes 29 of 30 evaluation questions; 172 of 193 numbers in its answers (89.1\%) are grounded in tool output.
\end{abstract}

\tableofcontents
\newpage

% ===========================================================================
\section{Overview}
% ===========================================================================

\begin{tabularx}{\linewidth}{@{}lL@{}}
\toprule
\textbf{App} & \textbf{Turnout}: one place to measure, diagnose, decide and manage attendee growth, with an AI analyst whose every number can be traced to the data. \\
\textbf{Challenge} & A: Attendee Growth~\cite{risa2025assignment} \\
\textbf{City} & Bengaluru (reasons below) \\
\textbf{Live link} & \LiveURL \\
\textbf{Passcode} & \Passcode \\
\textbf{Demo video} & \VideoURL \\
\textbf{3-slide deck} & \SlidesURL \\
\textbf{Source code} & \RepoURL \\
\bottomrule
\end{tabularx}

\paragraph{Why Bengaluru.}
\begin{itemize}
  \item \textbf{Largest addressable audience.} Bengaluru Urban has 1,148 colleges, the most of any district in India; Pune has 517 and Hyderabad 465~\cite{aishe2024}. Karnataka has the highest college density in the country~\cite{deccanherald2024}, and Bangalore University alone affiliates 500+ colleges with 3,00,000+ students~\cite{wikiBangaloreEdu}.
  \item \textbf{The channels in the dataset fit the city.} Ambassadors, college partnerships, club sign-ups and referrals all depend on many dense campuses close together.
  \item \textbf{A sourced number to plan from.} The college count lets the goal plan (\S\ref{sec:calc-plan}) start from a population instead of a guess.
\end{itemize}
Pune was the runner-up. Benchmark fests elsewhere show the range of scale (Mood Indigo 1,50,000+~\cite{wikiMoodIndigo}, Saarang 80,000+~\cite{iitmSaarang}, Oasis 8,000--10,000~\cite{wikiOasis}); CampusFest's goal of 2,200 verified visitors is a mid-size, single-venue event.

\subsection{60-second reviewer tour}
\begin{enumerate}
  \item Open the live link and enter the passcode. The welcome screen shows the Turnout mark; press anywhere. A nine-step tour follows on the first visit (\emph{Skip} is on every step; \emph{Help} at the top right replays it).
  \item On \textbf{Overview}, pick the \emph{Decision point} preset (31 Mar 2025). The headline reads \emph{``On pace for 1,837 visitors, 363 short of 2,200''}.
  \item Click the \textbf{Registrations} card to open its terminal-style drill-down: chart with goal line and forecast cone, by-campaign breakdown.
  \item Open \textbf{Funnel} (bottleneck highlighted), then \textbf{Bets} (five ranked strategies judged by early signals).
  \item Open \textbf{Copilot terminal} and ask, in order: \emph{``Are we on track?''}, \emph{``Why?''}, \emph{``What should we do?''}, \emph{``Create a task for that''}, and \emph{``What was our Instagram CTR in 2023?''}. Click any step in the trace to see its inputs.
  \item Press \textbf{$\langle$Cmd/Ctrl$\rangle$+J} and say \emph{``take me to paid social''} (needs a microphone).
\end{enumerate}

\shot{overview.png}{Overview as of the 31 March decision point: headline in words, the three goal metrics with projection, supporting signals, this week's drivers, and ranked actions.}{fig:overview}

% ===========================================================================
\section{Demo workflow}\label{sec:demo}
% ===========================================================================
The assignment asks for one complete workflow. All screenshots below are from the running app, as of 31 March 2025, in a single copilot conversation in the terminal. The left pane is the chat; the right pane is the live execution trace (model $\to$ tools $\to$ engine $\to$ data and rules $\to$ answer) with a verdict panel and an audit tape.

\begin{longtable}{@{}p{0.4cm}p{3.6cm}p{11.4cm}@{}}
\toprule
\textbf{\#} & \textbf{Ask or do} & \textbf{What happens} \\
\midrule
1 & \emph{Are we on track?} & \code{get\_overview}. Registrations \textbf{at risk}: 2,337 now, 2,887 projected against 3,200. Visitors \textbf{off track}: 1,837 projected against 2,200. Spend under budget. \\
2 & \emph{Why?} (follow-up) & \code{get\_bottleneck}, \code{get\_funnel}. Identified$\to$registered is 70.4\% against the plan-implied 88.9\% (index 0.79). Weakest campaigns: Creator Partners 61.0\%, Community Groups 67.4\%, Paid Social 68.2\%. \\
3 & \emph{What should we do?} & \code{get\_rule\_evaluations}, \code{suggest\_reallocation}. \textbf{Stop} Paid Social (\rs256 per registration over 14 days) and Creator Partners (\rs134); both are past the stop line of \rs110. Move \rs16,800 to the four cheapest scaling campaigns for +157 registrations and +100 visitors. 14 of 14 numbers grounded. \\
4 & \emph{Create a task for that} & \code{propose\_task} returns a draft only. A card shows title, owner, due date and priority, with \emph{Confirm} and \emph{Discard}. Nothing is written until \emph{Confirm}; the task then appears in Tasks and in the Activity log as created by \code{copilot}. \\
5 & Edit the data, ask again & On Campaigns $\to$ Community Groups $\to$ Data, log a \rs5,000 spend row for 31 Mar. Its 14-day cost per registration moves from \rs25 to \rs103.1 and the portfolio CPR from \rs73.58 to \rs80.87. The campaign's verdict and the copilot's answer change on the next request (evaluation case X1, \S\ref{sec:testing}). \emph{Reset} in Targets \& constraints restores the original data. \\
6 & \emph{What was our Instagram CTR in 2023?} & \code{get\_data\_coverage}. ``I don't have any 2023 data in the current dataset'', with the actual coverage (15 Jan--20 Apr 2025) and what data would be needed. \\
\bottomrule
\end{longtable}

\shot{demo-1-on-track.png}{Step 1: ``Are we on track?''. The answer names each goal's status; the trace shows the tool and engine functions that produced it.}{fig:demo1}
\shot{demo-2-why.png}{Step 2: the follow-up ``Why?'' inherits the context and goes to the bottleneck.}{fig:demo2}
\shot{demo-3-recommend.png}{Step 3: recommendation with the rule verdicts and the reallocation. Verdict panel: GROUNDED, 14/14 numbers. The audit tape records three Groq rate-limit fallbacks during this answer.}{fig:demo3}
\shot{demo-4-task.png}{Step 4: the task is proposed, not created. The card waits for Confirm or Discard.}{fig:demo4}
\shot{demo-5-campaign-c07.png}{Step 5: the Community Groups campaign page, where a daily data row is logged. The rule verdict on this page and the copilot both read the new data on the next request.}{fig:demo5}
\shot{demo-6-unanswerable.png}{Step 6: an unanswerable question gets an honest reply grounded in the data coverage, not a guess.}{fig:demo6}

% ===========================================================================
\section{Research and data sources}\label{sec:research}
% ===========================================================================
Research was done before building (full notes, with conflicts between sources, in \code{docs/research.md}). It covered fest attendance, the city's student base, ad costs, funnel and no-show benchmarks, and the free-tier limits of every service used. \textbf{Sponsor deal sizes were not researched}: they belong to Challenge B, which we did not choose.

\begin{longtable}{@{}p{4.3cm}p{4.0cm}p{1.2cm}p{5.6cm}@{}}
\toprule
\textbf{Metric} & \textbf{Value} & \textbf{Source} & \textbf{Label and use} \\
\midrule
\endhead
Colleges, Bengaluru Urban & 1,148 (most in India) & \cite{aishe2024} & Sourced. City choice; segment populations \\
Colleges, Pune / Hyderabad & 517 / 465 & \cite{aishe2024} & Sourced. Runner-up comparison \\
Karnataka higher-education enrolment & about 29.9 lakh & \cite{deccanherald2024} & Sourced. Context \\
Bangalore University & 3,00,000+ students, 500+ colleges & \cite{wikiBangaloreEdu} & Sourced (Wikipedia) \\
Bengaluru college students & about 8--10 lakh & derived & \textbf{Assumption}, derived from the two rows above. Goal plan uses 8,30,000 across five segments \\
Mood Indigo, IIT Bombay & 1,50,000+ & \cite{wikiMoodIndigo} & Sourced. Scale benchmark \\
Mood Indigo 2025 & Crowd crush; pre-registrations cancelled & \cite{careers360MoodIndigo} & Sourced. Why projections are capped at venue capacity \\
Saarang, IIT Madras & 80,000+ expected & \cite{iitmSaarang} & Sourced. Scale benchmark \\
Techfest, IIT Bombay & 1,65,000 (2015--16) & \cite{wikiTechfest} & Sourced. Scale benchmark \\
Oasis, BITS Pilani & 8,000--10,000 & \cite{wikiOasis} & Sourced. Closest in scale \\
Meta click-through rate, India & 1.0--1.5\% & \cite{owlclawMeta} & Sourced (vendor). CTR benchmark 1.0\% \\
Meta cost per click, India & \rs8--25 & \cite{owlclawMeta} & Sourced (vendor). Spend sanity check \\
Meta cost per lead, India & \rs150--400 & \cite{owlclawMeta} & Sourced (vendor). Upper bound for cost per registration \\
Instagram CPM, India & \rs45--350 & \cite{upgrowthInsta} & Sourced (vendor). Spend sanity check \\
Free-event no-show (median) & about 28\% (n=389) & \cite{pheedloop5} & Sourced. Base show-up \textbf{72\%} \\
Large free-event no-show & about 24\% & \cite{pheedloop5} & Sourced. Stretch show-up \textbf{76\%} \\
Paid-event no-show & about 17\% (n=683) & \cite{pheedloop5} & Sourced. Not used (CampusFest is free) \\
Free-event no-show, vendor rule & 40--60\% & \cite{nunify} & Sourced (vendor). Pessimistic bound for the what-if planner \\
Click $\to$ registration & 15--30\% & none found & \textbf{Assumption}. Session$\to$start benchmark 15\% \\
Click $\to$ session & 70\% & none found & \textbf{Assumption}. Benchmark \\
Registered $\to$ confirmed & 90\% & none found & \textbf{Assumption}. Benchmark \\
Sponsor deal sizes & not researched & n/a & Out of scope (Challenge B) \\
\bottomrule
\end{longtable}

\paragraph{The supplied dataset.} The RISA workbook~\cite{risa2025assignment} is event-sourced and person-level: 11 sheets, including 483 daily marketing rows, 3,685 people, 11,650 interactions, 2,840 registrations, 5,911 status versions and 3,390 check-in attempts, from 15 Jan to 20 Apr 2025. It is kept byte-for-byte in \code{data/original/} and the seed checks its SHA-256 before importing. Its full profile, including eleven deliberate traps (duplicates, internal and test people, unverified check-ins, status flips, vendor credits), is in \code{docs/dataset.md}.

% ===========================================================================
\section{Assumptions}\label{sec:assumptions}
% ===========================================================================
\subsection{Three kinds of data, kept apart}
\begin{tabularx}{\linewidth}{@{}p{2.8cm}LL@{}}
\toprule
\textbf{Kind} & \textbf{What it is} & \textbf{How it is marked in the app} \\
\midrule
\textbf{Original dataset} & The RISA workbook, unchanged. Every metric, verdict and projection is computed from it, after the dataset rules in \S\ref{sec:calc-rules}. & No badge. \emph{Reset} restores it exactly. \\
\textbf{Simulated} & Content we drafted so the strategy screens have something to judge: the goal plan, five bets, decision log entries, four starter tasks. None of it changes any dataset number. & Badged \lbl{SIMULATED}; labelled \lbl{Assumption} with its reasoning. \\
\textbf{User / copilot} & Anything created or edited in the app after seeding. & Badged \lbl{USER} or \lbl{COPILOT}; edited rows flagged; every change in the Activity log. \\
\textbf{Assumptions} & Parameters the data cannot give us (below). & Labelled \lbl{Assumption} or \lbl{Sourced} with a link wherever shown, including in copilot traces. \\
\bottomrule
\end{tabularx}

\subsection{Assumptions, reasons and sensitivity}
Sensitivities are computed by the engine's scenario function for 31 March (base projection: 1,837 visitors) and reproduced by \code{scripts/paper-numbers.mts}.

\begin{longtable}{@{}p{3.1cm}p{2.3cm}p{4.9cm}p{4.9cm}@{}}
\toprule
\textbf{Assumption} & \textbf{Value} & \textbf{Why it is reasonable} & \textbf{Sensitivity} \\
\midrule
\endhead
Show-up before the event & 72\% base, 76\% stretch & Sourced medians for free events~\cite{pheedloop5}. The data itself, measured after the event, gave 77.0\%. & The largest lever: each percentage point is about 25.5 visitors. 72\%$\to$60\%: 1,837$\to$1,531 ($-306$). 50\%$\to$40\%: 1,276$\to$1,021 ($-255$). \\
Walk-ins & 5\% of visitors & No source; the data's actual result was 125 of 2,050 (6.1\%). & 0\%: 1,745 ($-92$). 10\%: 1,939 ($+102$). \\
Venue capacity & 2,500 (editable) & The dataset has no capacity field; Mood Indigo 2025 shows why a cap matters~\cite{careers360MoodIndigo}. & Not binding at base. Binds the stretch plan (2,988$\to$2,500). At 1,800 it would cap the base case. \\
Registration close & End of the day before the event (11 Apr) & Last observed registration: 10 Apr. & Sets days left (11 on 31 Mar); each extra day adds about 50 registrations at the current pace. \\
Extra-spend efficiency & Added budget converts at 2$\times$ the recipient's recent CPR & Diminishing returns: the cheapest audience is reached first. & A less conservative 1$\times$ would show a larger gain from the same move; we report the conservative figure. \\
Rule thresholds & Scale $\le$0.8$\times$, stop $\ge$1.5$\times$ portfolio CPR; $\ge$15 registrations & The source has no per-campaign cost targets, so the portfolio is the yardstick. & On 31 Mar: scale $\le$\rs58.86, stop $\ge$\rs110.37. Paid Social (\rs256) and Creator Partners (\rs134) stop under any threshold from 1.2$\times$ to 1.8$\times$. \\
Funnel benchmarks & CTR 1.0\%; click$\to$session 70\%; session$\to$start 15\%; confirm 90\% & CTR sourced~\cite{owlclawMeta}; others assumed. & Only identified$\to$registered (index 0.79) and confirm (0.93) fall below 1; the bottleneck does not change unless the confirm benchmark rises above about 106\%. \\
Goal-plan segment rates & Reach 3--25\%, identify 8--12\%, register 75--85\% & Planning assumptions per segment (\S\ref{sec:calc-plan}). & Stretch multiplies reach by 1.25 and registration by 1.08: 2,097$\to$2,988 visitors before the cap. \\
Bengaluru student population & 8,30,000 across five segments & Derived from AISHE college counts~\cite{aishe2024}. & Linear: 10\% fewer students gives 10\% fewer planned visitors. \\
\bottomrule
\end{longtable}

% ===========================================================================
\section{Detailed calculations}\label{sec:calc}
% ===========================================================================
Every number comes from one engine (\code{lib/analytics}), made of pure functions that each return \code{\{value, formula, inputs, assumptions, sources\}}. The figures below are asserted in \code{tests/engine.golden.test.ts} against an independent pandas computation, and \code{npm run report -{}- 2025-03-31} prints them. The longer sheet is \code{docs/calculations.md}.

\subsection{Dataset rules applied before any metric}\label{sec:calc-rules}
\begin{tabularx}{\linewidth}{@{}p{4.2cm}L@{}}
\toprule
\textbf{Rule} & \textbf{Effect on the original data} \\
\midrule
As of = end of that IST day & No later record is counted (no look-ahead). \\
Deduplicate by source ID & Removes 2 marketing rows, 4 status versions, 20 check-ins. Impressions: 1,032,918 summed, 1,029,015 deduplicated. \\
External people only & Excludes 40 internal, 10 staff and 10 test identities. \\
Latest status at the cutoff & Cancel $\to$ reconfirm handled correctly. \\
Verified attendance & Admitted \emph{and} verified, external, distinct person. Excludes 50 denied and 30 pending check-ins; repeat attempts count once. \\
Target in force & Launch plan until 15 Mar, then the revised plan (3,600 identified, 3,200 registrations, 2,200 visitors, \rs3,00,000). \\
Net spend & Vendor credits subtract (\rs2,500). \\
\bottomrule
\end{tabularx}

\subsection{Funnel (two layers, never divided across)}
Anonymous events (impressions $\to$ clicks $\to$ sessions $\to$ form starts) and identified people (identified $\to$ engaged $\to$ registered $\to$ confirmed $\to$ attended) are separate layers; a conversion is a stage divided by the stage before it in the same layer. Final state on 20 April:

\begin{center}\small
\begin{tabular}{@{}lrr@{\hspace{2.4em}}lrr@{}}
\toprule
\textbf{Anonymous} & \textbf{Count} & \textbf{Conv.} & \textbf{Identified} & \textbf{Count} & \textbf{Conv.} \\
\midrule
Impressions & 1,029,015 & & Identified (campaign) & 3,500 & \\
Clicks & 32,070 & 3.12\% & Engaged & 3,317 & 94.8\% \\
Sessions & 24,645 & 76.8\% & Registered & 2,800 & 84.4\% \\
Form starts & 5,548 & 22.5\% & Confirmed & 2,500 & 89.3\% \\
 & & & Attended (registered) & 1,925 & 77.0\% \\
 & & & + walk-ins & 125 & \\
 & & & \textbf{Unique verified visitors} & \textbf{2,050} & \\
\bottomrule
\end{tabular}
\end{center}

Final result against the revised plan: identified 3,625 / 3,600 (\textbf{achieved}); registrations 2,800 / 3,200 = 87.5\% (\textbf{missed}); visitors 2,050 / 2,200 = 93.2\% (\textbf{missed}); spend \rs2,74,000 / \rs3,00,000 (under budget).

\subsection{The goal plan: base and stretch, bottom up}\label{sec:calc-plan}
The plan answers ``how high can we aim?'' from the audience rather than from last year. For each segment $s$:
\[
\text{identified}_s = P_s \cdot r_s \cdot i_s,\qquad
\text{registered}_s = \text{identified}_s \cdot g_s
\]
where $P$ is population, $r$ reach, $i$ identify rate and $g$ register rate. Then, with confirm rate $c$, show-up $u$, walk-in share $w$ and capacity $K$:
\[
\text{visitors} = \min\!\left(K,\ \frac{\sum_s \text{registered}_s \cdot c \cdot u}{1-w}\right)
\]

\begin{center}\small
\begin{tabular}{@{}lrrrrrrr@{}}
\toprule
\textbf{Segment} & $P$ & $r$ & \textbf{Reached} & $i$ & \textbf{Identified} & $g$ & \textbf{Registered} \\
\midrule
Engineering \& tech & 2,40,000 & 5\% & 12,000 & 12\% & 1,440 & 80\% & 1,152.0 \\
Arts, science \& commerce & 4,20,000 & 3\% & 12,600 & 8\% & 1,008 & 75\% & 756.0 \\
Management \& design & 90,000 & 6\% & 5,400 & 10\% & 540 & 78\% & 421.2 \\
Clubs \& societies & 60,000 & 10\% & 6,000 & 9\% & 540 & 85\% & 459.0 \\
Friends of registrants & 20,000 & 25\% & 5,000 & 8\% & 400 & 80\% & 320.0 \\
\midrule
\textbf{Total} & 8,30,000 & & 41,000 & & \textbf{3,928} & & \textbf{3,108.2} \\
\bottomrule
\end{tabular}
\end{center}

\noindent\textbf{Worked example (base).} Engineering: $2{,}40{,}000 \times 0.05 = 12{,}000$ reached; $\times\,0.12 = 1{,}440$ identified; $\times\,0.80 = 1{,}152$ registered. Summing all segments:
\begin{align*}
\text{confirmed} &= 3{,}108.2 \times 0.89 = 2{,}766.3 &&(c = 89\%, \text{ derived from the data})\\
\text{registered visitors} &= 2{,}766.3 \times 0.72 = 1{,}991.7 &&(u = 72\%, \text{ sourced})\\
\text{visitors} &= 1{,}991.7 \div 0.95 = \mathbf{2{,}096.6} &&(w = 5\%;\ 104.8 \text{ walk-ins}; \text{ below } K = 2{,}500)
\end{align*}
\textbf{Stretch} (reach $\times1.25$, registration $\times1.08$, show-up 76\%): identified $3{,}928 \times 1.25 = 4{,}910$; registered $3{,}108.2 \times 1.25 \times 1.08 = 4{,}196.1$; confirmed $\times0.89 = 3{,}734.5$; registered visitors $\times0.76 = 2{,}838.2$; visitors $2{,}838.2 \div 0.95 = 2{,}987.6$, \textbf{capped at 2,500}.

So the revised goal of 2,200 sits between base (2,097) and stretch (2,500): reachable, but only with better reach or conversion than the base case. The app compares each plan step with the data (plan vs data): on 31 March identified is 3,320 against 3,928 planned, and identified$\to$registered is 70.4\% against the plan's 79.1\%.

\diagram{funnel-math.jpg}{Base and stretch goal plans, from audience to visitors, with the capacity cap.}{fig:funnelmath}

\subsection{Projection from the data (what the dashboard shows)}
\[
\text{reg}_{\text{proj}} = \sum_c \big(\text{reg}_c + \text{rate}_c \cdot d\big),\qquad
\text{visitors} = \min\!\left(K,\ \frac{\text{reg}_{\text{proj}} \cdot \text{confirmed share} \cdot u}{1-w}\right)
\]
with $d$ the registration days left. \textbf{Base}: $\text{rate}_c$ = the campaign's trailing 7-day registrations $\div$ 7, $u = 72\%$. \textbf{Stretch}: its best 7-day window so far, $u = 76\%$.

\noindent\textbf{Worked example, 31 March (base).}
\begin{enumerate}
  \item Trailing 7 days (25--31 Mar): 350 registrations, so 50.0 per day.
  \item Registrations $= 2{,}337 + 50 \times 11 = \mathbf{2{,}887}$.
  \item Confirmed share $= 1{,}962 \div 2{,}337 = 83.95\%$, so $2{,}887 \times 0.8395 = 2{,}423.8$ confirmed.
  \item Registered visitors $= 2{,}423.8 \times 0.72 = 1{,}745.1$; visitors $= 1{,}745.1 \div 0.95 = \mathbf{1{,}837}$.
  \item Spend $= 2{,}48{,}150 + (26{,}500 \div 7) \times 11 = \rs\mathbf{2{,}89{,}793}$.
\end{enumerate}
Stretch on 31 March: 3,077 registrations and 2,067 visitors. Even the stretch case misses 2,200 unless conversion improves: that is the decision point.

\subsection{Target versus actual}
$\text{pct} = \text{actual} \div \text{goal}$; required per day $= (\text{goal} - \text{actual}) \div d$; current per day = trailing 7-day average. Status from the base projection: $\ge$ goal \emph{on track}, $\ge 90\%$ \emph{at risk}, otherwise \emph{off track}; after the window closes, \emph{achieved} or \emph{missed}. On 31 March, registrations need $(3{,}200 - 2{,}337) \div 11 = 78.5$ per day against a current 50.0, and $2{,}887 / 3{,}200 = 90.2\%$, hence \emph{at risk}.

\subsection{Bottleneck}
For each stage, index $=$ actual conversion $\div$ benchmark; the bottleneck is the lowest index below 1 with at least 50 records. On 31 March:

\begin{center}\small
\begin{tabular}{@{}lrrrl@{}}
\toprule
\textbf{Stage} & \textbf{Actual} & \textbf{Benchmark} & \textbf{Index} & \textbf{Weakest campaigns} \\
\midrule
Click-through & 3.1\% & 1.0\% & 3.12 & Paid Social 2.2\% \\
Click $\to$ session & 76.7\% & 70\% & 1.10 & Student Referrals 73.9\% \\
Session $\to$ start & 22.9\% & 15\% & 1.53 & College Partnerships 15.9\% \\
\textbf{Identified $\to$ registered} & \textbf{70.4\%} & \textbf{88.9\%} & \textbf{0.79} & Creator Partners 61.0\%, Community Groups 67.4\% \\
Registered $\to$ confirmed & 84.0\% & 90\% & 0.93 & Creator Partners 77.7\% \\
\bottomrule
\end{tabular}
\end{center}
The identified$\to$registered benchmark is derived from the plan itself: $3{,}200 \div 3{,}600 = 88.9\%$.

\subsection{Scale, keep, change and stop rules}
Each campaign is judged on its cost per registration (CPR) over the last 14 days, against the portfolio CPR for the same window:

\begin{center}\small
\begin{tabular}{@{}ll@{}}
\toprule
\textbf{Verdict} & \textbf{Condition} \\
\midrule
Insufficient data & fewer than 15 registrations in the window \\
Scale & CPR $\le 0.8 \times$ portfolio CPR \\
Stop & CPR $\ge 1.5 \times$ portfolio CPR (or already stopped) \\
Change & within thresholds, but CTR, click$\to$session or session$\to$start $<80\%$ of benchmark, \\ & cancellations $>15\%$, or weak show-up \\
Keep & otherwise \\
\bottomrule
\end{tabular}
\end{center}
\noindent On 31 March the portfolio CPR is $50{,}475 \div 686 = \rs73.58$, so scale at $\le\rs58.86$ and stop at $\ge\rs110.37$. Paid Social: $19{,}200 \div 75 = \rs256$ (\textbf{stop}). Creator Partners: $12{,}200 \div 91 = \rs134$ (\textbf{stop}). Community Groups \rs25, Student Referrals \rs30, Organic Search \rs33, Campus Ambassadors \rs35, College Partnerships \rs49 (all \textbf{scale}).

\subsection{Reallocation and what-if scenarios}
Future spend per campaign = trailing 7-day daily spend $\times$ days left. Stop campaigns give 100\% of it, change campaigns 50\%; scale campaigns receive it cheapest first, each capped at its own planned future spend. Removed money loses registrations at the donor's CPR; added money gains them at $2\times$ the recipient's CPR. On 31 March this moves \rs16,800 (\rs15,000 from Paid Social, \rs1,800 from Creator Partners) to Student Referrals, Organic Search, Campus Ambassadors and College Partnerships: \textbf{+157 registrations, +100 visitors}. The what-if planner runs the same projection with changes applied (shift or add budget, scale or stop a campaign, set show-up, walk-ins or capacity); for example show-up at 60\% gives $2{,}423.8 \times 0.60 \div 0.95 = 1{,}531$ visitors.

% ===========================================================================
\section{Prioritisation framework}\label{sec:rice}
% ===========================================================================
Two layers, because two different questions are being answered.

\paragraph{Planning: RICE~\cite{intercomRice}.} Which strategies deserve budget before there is evidence? $\text{RICE} = \text{Reach} \times \text{Impact} \times \text{Confidence} \div \text{Effort}$, with reach in students per campaign period, impact on a 0.25--3 scale, confidence 0--1 and effort in person-weeks.

\begin{center}\small
\begin{tabular}{@{}clrrrrrrr@{}}
\toprule
\textbf{\#} & \textbf{Bet} & \textbf{Reach} & \textbf{I} & \textbf{C} & \textbf{E} & \textbf{RICE} & \textbf{Budget} & \textbf{Visitors} \\
\midrule
1 & Campus ambassador network & 40,000 & 2 & 0.8 & 6 & \textbf{10,667} & \rs60,000 & 450 \\
2 & College club group sign-ups & 25,000 & 2 & 0.8 & 4 & \textbf{10,000} & \rs40,000 & 420 \\
3 & Bring-3-friends referrals & 15,000 & 2 & 0.7 & 3 & \textbf{7,000} & \rs25,000 & 300 \\
4 & Creator-led reels & 80,000 & 1 & 0.4 & 5 & \textbf{6,400} & \rs60,000 & 250 \\
5 & Paid social retargeting & 1,20,000 & 0.5 & 0.4 & 4 & \textbf{6,000} & \rs95,000 & 300 \\
\bottomrule
\end{tabular}
\end{center}
Worked example: bet 1 $= 40{,}000 \times 2 \times 0.8 \div 6 = 10{,}667$. Peer channels rank first because they combine high impact on the bottleneck stage (identified$\to$registered) with confidence; the two broad-reach media bets have the most reach but low confidence and low impact per person. Each bet carries early signals with scale and stop lines (for example CPR over 14 days), so the ranking is re-judged by the data as it arrives (Figure~\ref{fig:bets}).

\paragraph{Operating: rules.} Once campaigns are running, RICE estimates are replaced by evidence: the CPR rules above. The two frameworks agree here. The two lowest-RICE bets map to Creator Partners and Paid Social, which the rules later mark \emph{stop}.

\paragraph{Alternatives we dropped.}
\begin{itemize}
  \item \textbf{ICE} (impact, confidence, ease): no reach term, but reach is exactly what differs most between peer and media channels.
  \item \textbf{MoSCoW / Kano}: categorise but do not rank within a category, and do not use numbers the data can later check.
  \item \textbf{Weighted scoring with many criteria}: more knobs than evidence; harder to explain in one line.
  \item \textbf{Cost per visitor as the operating metric}: the right goal metric, but it only exists after the event. CPR is available every day and correlates with it.
\end{itemize}

\shot{bets.png}{Bets: five ranked strategies, each judged live by its early signals (scale / watch / stop / too early / from event day).}{fig:bets}

% ===========================================================================
\section{Product and architecture}\label{sec:arch}
% ===========================================================================
\diagram{strategy-loop.jpg}{The operating loop the app is built around. The as-of date drives every step; the copilot can help at each one, and a person confirms every change.}{fig:loop}
\begin{tabularx}{\linewidth}{@{}p{2.2cm}p{4.2cm}L@{}}
\toprule
\textbf{Module} & \textbf{Screens} & \textbf{What it does} \\
\midrule
\textbf{Measure} & Overview, KPI terminals (REG, VIS, IDN, SPD, CPR, SHW, CTR, CNF), Funnel, Campaigns & Goals vs actual and projection; two-layer funnel; every KPI drills down to a terminal-style chart with goal line, forecast cone and by-campaign breakdown. A global \emph{as of} date replays any day from 15 Jan to 20 Apr 2025. \\
\textbf{Diagnose} & Funnel, Overview (``This week''), Campaign detail & Bottleneck stage and weakest campaigns; weekly drivers (association, not cause); per-campaign health checks. \\
\textbf{Decide} & Campaigns, What-if planner, Goal plan, Bets & Rule verdicts; budget reallocation; scenarios; bottom-up plan; ranked bets with early signals. \\
\textbf{Manage} & Tasks, Decisions, Team, Activity, Targets \& constraints & Tasks with owners and due dates; a decision log whose outcome is measured 14 days later; workload per owner; a full change log; editable constraints; reset. \\
\textbf{Ask} & Copilot terminal, chat drawer, voice & Grounded answers, traces, proposals that wait for a yes, navigation by voice. \\
\bottomrule
\end{tabularx}

\diagram{architecture.jpg}{System architecture. Pages, the REST API and the copilot all call the same pure analytics engine, so the chatbot cannot disagree with the screen.}{fig:arch}

\paragraph{Collections (MongoDB Atlas).} Source tables mirrored one-to-one with source IDs as unique keys: \code{campaigns}, \code{campaign\_daily}, \code{people}, \code{interactions}, \code{registrations}, \code{registration\_history}, \code{checkins}, \code{targets}, \code{event}, each with a \code{baseline\_*} twin used by reset. App data: \code{tasks}, \code{bets}, \code{decisions}, \code{goal\_model}, \code{settings}, \code{activity}, \code{chat\_sessions}, and \code{meta} (holding \code{data\_version}).

\paragraph{How a change reaches the analytics.} All writes go through one path: validate with Zod, write, append to the activity log, then increment \code{meta.data\_version}. Each request reads only that counter; when it has changed, the dataset is reloaded (memory $\to$ compressed disk pack $\to$ MongoDB) and every view, KPI and copilot tool recomputes from it. On the free M0 cluster this cut a cold start from about 2.5 minutes to about 1.5 seconds, with warm pages in 0.06--0.8 seconds.

\diagram{data-flow.jpg}{From an edit to every answer: one write path, one version counter, one engine.}{fig:dataflow}

% ===========================================================================
\section{AI chatbot design}\label{sec:copilot}
% ===========================================================================
\diagram{copilot-loop.jpg}{How the copilot answers: tool router, provider chain, tools, engine, data and rules, and the claim check.}{fig:copilotloop}

\paragraph{Grounded in the data.} The model receives a short header (as-of date, phase, target, campaign list) and 24 tools that call the engine: overview, funnel, bottleneck, trends, daily series, campaigns (with fuzzy name matching), comparisons, rule evaluations, reallocation, scenarios, tasks, activity, data coverage, goal plan, bets, weekly drivers, decisions and team. A \code{calculate} tool (a sandboxed mathjs) does all arithmetic; the model is told never to compute in its head. After the answer, every number in it is matched against the numbers returned by tools in this turn or earlier turns; unmatched numbers are flagged \textbf{UNSUPPORTED} in the trace, and the verdict panel reads GROUNDED or CHECK.

\paragraph{Confirm before writing.} The three write tools (\code{propose\_task}, \code{propose\_campaign\_update}, \code{propose\_decision}) only validate and return a draft. The interface shows a card (Figure~\ref{fig:demo4}); \emph{Confirm} calls the same write path as the screens with \code{actor: "copilot"}, so the change is logged and badged. The voice assistant follows the same rule: it asks ``Shall I\ldots?'' and acts only on a clear yes, or on the Accept button.

\paragraph{Uncertainty and follow-ups.} Before saying the data cannot answer, the copilot must call \code{get\_data\_coverage}; it then states the coverage and what data would be needed (Figure~\ref{fig:demo6}). Answers separate evidence, assumptions (labelled sourced or assumed) and confidence, and distinguish association from cause. Follow-ups (``Why?'', ``And Creator Partners?'') work because the tool router looks at the last two user messages, and earlier turns are kept as text; ambiguous names (``the partners campaign'') return candidates instead of a guess.

\paragraph{Provider chain.} Groq first: \code{gpt-oss-120b} $\to$ \code{qwen3.8-27b} $\to$ \code{gpt-oss-20b}, across two keys, then Gemini \code{3.8-flash} $\to$ \code{3.5-flash}~\cite{groqLimits,geminiModels}. It is implemented as model middleware, so fallback happens per step inside a multi-step tool loop, and every attempt is shown in the trace. To stay within Groq's free tier (about 8k tokens per minute), a router offers only the tools relevant to the question (at most 10).

\paragraph{Execution trace.} Recorded on the server from what actually ran, never from the model describing itself, and streamed to the browser. It is drawn as stage columns (Model $\to$ Tools $\to$ Engine $\to$ Data \& rules $\to$ Answer); clicking a step shows its inputs, output, formula and sources next to it.

\paragraph{Voice assistant.} A bottom-centre mark starts a Gemini Live session~\cite{geminiLive}. The browser gets a single-use, short-lived token from the server, so the key never leaves it~\cite{geminiEphemeral}. Its tools are the same engine tools (run on the server), the same propose-then-confirm write tools, and screen tools run in the browser: navigate (``take me to paid social''), offer navigation for ``where is\ldots'' questions, read the current screen, go back, and show long answers as an on-screen card. It was tested with real speech and in a browser test with a simulated microphone (five of five scripted turns succeeded); spoken replies start within 1--3 seconds (Figure~\ref{fig:voice}).

\shot{voice.png}{Voice assistant: a proposed task waiting for a yes; the bar shows the spoken reply as a caption.}{fig:voice}

\paragraph{Limits.} Free-tier rate limits cause fallbacks and occasional queueing (visible in Figures~\ref{fig:demo3}--\ref{fig:demo6}); a provider error after the first token is not retried on another provider; claim matching is numeric, so a coincidentally equal small number can count as grounded; multi-step chat answers take about 5--13 seconds on free-tier models (voice replies in 1--3 seconds); and the copilot can only see what the dataset and the app contain.

% ===========================================================================
\section{Testing}\label{sec:testing}
% ===========================================================================
\paragraph{Unit and golden tests (\code{npm test}, 60 passing).} Hand-built fixtures for edge cases: zero denominators, an empty campaign, dates before the history, look-ahead, cancel $\to$ reconfirm, duplicates, the capacity cap and scenario arithmetic. Golden tests on the real workbook, against an independent pandas computation on three dates. The strategy engine (goal-plan arithmetic checked by hand, bet signals on real data, decision outcomes, weekly drivers, owner workload). Claim extraction and grounding, and tool routing.

\paragraph{Chatbot evaluation (\code{npm run eval}).} 30 questions across 11 categories, run against the live app. Expected values are computed by the engine at run time, not hard-coded, so the test stays valid after data changes. \textbf{29 of 30 passed (96.7\%)}; 172 of 193 numbers in answers (89.1\%) were grounded in tool output. The first run scored 23 of 30; its seven failures (a tool not offered, malformed tool arguments, missing walk-in and show-up fields, an ambiguous metric name) led to four fixes. The remaining failure is ambiguity: ``the partners campaign'' was answered for one campaign instead of asking which. The full table is in Appendix~\ref{app:eval} and \code{docs/eval-results.md}.

\begin{center}\small
\begin{tabular}{@{}lr@{\hspace{3em}}lr@{}}
\toprule
\textbf{Category} & \textbf{Passed} & \textbf{Category} & \textbf{Passed} \\
\midrule
Lookup & 5/5 & Action (proposes, never writes) & 2/2 \\
Calculation & 4/4 & Unanswerable & 4/4 \\
Comparison & 3/3 & Ambiguous & 0/1 \\
Diagnosis & 3/3 & Data quality & 2/2 \\
Recommendation & 3/3 & After a data change & 1/1 \\
Follow-up & 2/2 & \textbf{Total} & \textbf{29/30} \\
\bottomrule
\end{tabular}
\end{center}

\paragraph{Browser tests (real Chrome).} \code{scripts/onboarding-e2e.mjs}: welcome $\to$ overview $\to$ all nine tour steps, no replay after finishing, Help replays, Skip closes, deep links skip the welcome. \code{scripts/assistant/voice-e2e.mjs}: voice with a simulated microphone and typed turns, covering navigation by name, a data question, a proposal and a spoken ``no'' (5/5). \code{scripts/paper-shots.mjs} produced every screenshot in this document.

\paragraph{Checked by hand.} The worked examples in \S\ref{sec:calc} against pandas; the data-change loop (a logged row moves Community Groups' CPR from \rs25 to \rs103.1 and the portfolio CPR from \rs73.58 to \rs80.87, and the copilot's answer follows); every screen at the decision point, final day and launch; fallbacks between providers under rate limits (visible in the traces); dark and light themes; reset to the original data.

% ===========================================================================
\section{Tools used}\label{sec:tools}
% ===========================================================================
\begin{tabularx}{\linewidth}{@{}p{3.6cm}p{4.2cm}L@{}}
\toprule
\textbf{Tool} & \textbf{Plan} & \textbf{Used for} \\
\midrule
Next.js 16.4, React 19, TypeScript & Open source & App framework~\cite{nextjs} \\
Vercel & Hobby (free) & Hosting~\cite{vercelLimits,vercelPool} \\
MongoDB Atlas & M0 (free) & Database~\cite{atlasM0} \\
Groq & Free tier, two keys & Primary copilot models~\cite{groqLimits} \\
Google AI Studio (Gemini) & Free tier & Fallback models; Gemini Live voice~\cite{geminiRate,geminiLive} \\
Vercel AI SDK 7 & Open source & Tool calling, streaming~\cite{aisdk7} \\
React Flow, dagre & Open source (MIT) & Trace graph~\cite{reactflow} \\
React Bits & MIT + Commons Clause & 13 animated components, restyled~\cite{reactbits} \\
Doto, Geist & SIL OFL & Dot-matrix display and UI fonts~\cite{doto} \\
Zod, mathjs, Tailwind 4 & Open source & Validation, calculator, styling \\
Vitest, Playwright & Open source & Tests, browser checks, screenshots \\
\bottomrule
\end{tabularx}

\paragraph{Reused, modified and verified.}
\begin{itemize}
  \item \textbf{Reused:} the Next.js scaffold; React Bits components (Counter, CountUp, DecryptedText, DotGrid, SplitFlapText, LatticeLoader, StatusMark, FaultyTerminal, Scanner, Noise, GlitchText, LetterGlitch, TextType); React Flow; dagre; Doto and Geist fonts.
  \item \textbf{Modified:} every React Bits component, restyled to monochrome and given a static fallback under reduced motion.
  \item \textbf{Written for this project:} the analytics engine, the write path and views, the copilot (tools, router, provider chain, trace recorder, claim grounding), the voice assistant, all pages, the dot-matrix chart components, the onboarding, and the test and evaluation scripts.
  \item \textbf{Verified:} engine figures against pandas; the worked examples here; the eval with expected values computed at run time; the data-change loop end to end; voice and onboarding in a real browser.
\end{itemize}

% ===========================================================================
\section{Setup}\label{sec:setup}
% ===========================================================================
\begin{verbatim}
git clone <repo> && cd turnout
npm install                       # Node 22 or newer
cp .env.example .env              # fill in the variables below
npm run seed                      # import data/original (checks SHA-256)
npm run seed:strategy             # goal plan, bets, decisions (simulated)
npm run verify-seed               # headline numbers straight from MongoDB
npm run dev                       # http://localhost:3000
\end{verbatim}

{\small\begin{tabularx}{\linewidth}{@{}p{6.4cm}L@{}}
\toprule
\textbf{Variable} & \textbf{Purpose} \\
\midrule
\code{MONGODB\_URI}, \code{MONGODB\_DB} & Atlas connection; database name (default \code{turnout}) \\
\code{GROQ\_API\_KEY}, \code{GROQ\_API\_KEY2} & Primary models; a second key doubles the free budget \\
\code{GOOGLE\_GENERATIVE\_AI\_API\_KEY}, \newline \code{GOOGLE\_API\_KEY\_2} & Gemini fallback \\
\code{ASSISTANT\_GEMINI\_API\_KEY} & Optional separate key for voice \\
\code{LLM\_PROVIDERS} & Order of providers (default \code{groq,google}) \\
\code{ADMIN\_TOKEN} & Protects seed and reset endpoints \\
\code{APP\_PASSCODE} & Optional passcode gate for the demo \\
\bottomrule
\end{tabularx}}

\paragraph{Reset to the original data.} Targets \& constraints $\to$ Reset; or \code{npm run seed -{}- -{}-reset}; or \code{POST /api/admin/reset} with \code{Authorization: Bearer \$ADMIN\_TOKEN}.

\paragraph{Deploy (GitHub $\to$ Vercel).} Push to GitHub and import the repository in the Vercel dashboard; add the variables above; set Node 22 and a Functions region near the Atlas cluster; allow \code{0.0.0.0/0} in Atlas network access (Hobby has no static IPs); on a fresh database call \code{POST /api/admin/seed} once with the bearer token.

% ===========================================================================
\section{Known limitations and next steps}\label{sec:limits}
% ===========================================================================
\begin{itemize}
  \item \textbf{Simulated strategy content.} The goal plan, bets, decision log and starter tasks were drafted by us to make the strategy screens meaningful; they are labelled as such and change no dataset number. Decision outcomes are associations over 14 days, not causal estimates.
  \item \textbf{Show-up before the event} uses a sourced benchmark, because the data cannot reveal it until the fest; it is also the largest sensitivity (\S\ref{sec:assumptions}).
  \item \textbf{Free-tier limits.} Groq's free tier causes fallbacks during bursts; a paid tier would remove them. A provider failure mid-stream is not retried.
  \item \textbf{Chat latency.} Multi-step chat answers take about 5--13 s on free-tier models; voice replies in 1--3 s.
  \item \textbf{Grounding is numeric.} It catches invented numbers, not wrong reasoning over correct numbers; coincidentally equal small numbers can count as grounded.
  \item \textbf{Single instance.} The rate limiter and the dataset cache are per instance; correctness does not depend on them because freshness comes from the version counter.
  \item \textbf{Not built:} Challenge B (sponsorship).
\end{itemize}
\textbf{Next:} a paid model tier or smaller task-specific prompts to bring answers under 3 s; persistent chat history; per-campaign budget plans; a cost-per-visitor rule once attendance data accrues; a larger evaluation set written by someone other than the builder.

% ===========================================================================
\section{Decisions I made}\label{sec:decisions}
% ===========================================================================
\begin{longtable}{@{}p{4.4cm}p{10.8cm}@{}}
\toprule
\textbf{Decision} & \textbf{Trade-off} \\
\midrule
\textbf{Challenge A over B} & The supplied dataset is a complete, person-level attendee history, so attendee growth could be built on real numbers end to end; sponsorship would have rested mostly on assumptions. Cost: no sponsor pipeline. \\
\textbf{Propose, never auto-execute} & A growth team acting on an AI suggestion should see exactly what will change and own the decision. Cost: one extra click per action. Every copilot change is badged and logged. \\
\textbf{One engine for screens, API and copilot} & The chatbot cannot give a number the dashboard does not show, and tests cover both at once. Cost: every new question type needs an engine function and a tool. \\
\textbf{Tools and a calculator instead of putting data in the prompt} & Smaller prompts that fit free-tier limits, exact arithmetic, and every number traceable. Cost: multi-step answers are slower. \\
\textbf{As-of replay instead of ``today''} & The dataset is a finished campaign; replaying any day lets the reviewer see the 31 March decision as it looked then, without look-ahead. Cost: every query carries a date. \\
\textbf{Portfolio-relative rule thresholds} & The source has no per-campaign cost targets; comparing with the portfolio adapts as costs move. Cost: in a uniformly bad portfolio nothing is marked stop, so custom thresholds can be set per campaign. \\
\bottomrule
\end{longtable}

% ===========================================================================
\section{Ownership and communication}\label{sec:ownership}
% ===========================================================================
\paragraph{How I decided.} Research came first and was written down with sources and conflicts (\code{docs/research.md}); every product decision was logged with its date and reason (\code{docs/decisions.md}). The dataset was profiled before any feature was designed (\code{docs/dataset.md}), which is how the eleven traps became engine rules with tests rather than surprises. Wherever the data could not settle a question, the app says so and labels the assumption.

\paragraph{How I tested.} Independent numbers first (pandas against TypeScript), then behaviour (fixtures for edge cases), then the AI (an evaluation whose expected answers are computed from live data), then the product in a real browser. Failures changed the product: the first evaluation's 23/30 led to four fixes; the browser test for voice found a race in turn-taking, which was fixed before it shipped.

\paragraph{What is still weak.}
\begin{itemize}
  \item The evaluation set was written by the builder; 30 questions are a smoke test, not a benchmark. The ambiguity case still fails.
  \item 89.1\% of numbers grounded means about one in ten numbers is either derived in text or unverifiable; they are flagged, not blocked.
  \item The strategy layer's inputs (segment rates, RICE scores, bet budgets) are our planning assumptions; they should be replaced with the team's own estimates.
  \item Multi-step chat answers are the slowest part of the app.
\end{itemize}

% ===========================================================================
\appendix
\section{Chatbot evaluation cases}\label{app:eval}
Run on 26 Nov 2026 against the running app. ``Grounded'' counts numbers in the answer matched to tool output.

{\small
\begin{longtable}{@{}p{0.6cm}p{2.3cm}p{7.2cm}p{0.9cm}p{3.4cm}@{}}
\toprule
\textbf{ID} & \textbf{Category} & \textbf{Question} & \textbf{Result} & \textbf{Check} \\
\midrule
\endhead
L1 & lookup & How many registrations do we have so far? & pass & 2,337 \\
L2 & lookup & What's our net marketing spend to date? & pass & \rs2,48,150 \\
L3 & lookup & How many unique verified visitors attended? & pass & 2,050 \\
L4 & lookup & How many walk-ins did we get? & pass & 125 \\
L5 & lookup & Which target version is in force and what's the visitor goal? & pass & revised, 2,200 \\
C1 & calculation & What percentage of the registration goal have we reached? & pass & 73.0\% \\
C2 & calculation & How many registrations per day do we need to hit the goal? & pass & 78.5 \\
C3 & calculation & What's our overall cost per registration to date? & pass & \rs106.2 \\
C4 & calculation & What was the show-up rate of confirmed registrants? & pass & 77\% \\
P1 & comparison & Compare Paid Social and College Partnerships on 14-day CPR. & pass & \rs256 vs \rs48.61 \\
P2 & comparison & Which campaign has brought the most registrations? & pass & Paid Social, 636 \\
P3 & comparison & Which campaign has the lowest click-through rate? & pass & Paid Social \\
D1 & diagnosis & Where's our biggest funnel bottleneck? & pass & identified$\to$registered 70.4\% \\
D2 & diagnosis & How did registrations change week over week? & pass & 350 vs 336 \\
D3 & diagnosis & Are we on track for the stretch target? & pass & 2,066.7 visitors \\
R1 & recommendation & Which campaign should we cut? & pass & Paid Social, rules \\
R2 & recommendation & Where should we move budget, and how much? & pass & \rs16,800 \\
R3 & recommendation & What if show-up drops to 60\%? & pass & 1,530.8 \\
F1 & follow-up & Paid Social's 14-day CPR? $\to$ And Creator Partners? & pass & \rs134.07 \\
F2 & follow-up & Lowest 14-day CPR? $\to$ Who owns it? & pass & Nisha Rao \\
A1 & action & Create a task for Arjun Verma to pause Paid Social, due 3 April. & pass & proposed, not created \\
A2 & action & Pause the Paid Social campaign. & pass & proposed update \\
U1 & unanswerable & What was our Instagram CTR in 2023? & pass & says the data can't answer \\
U2 & unanswerable & What was the average age of our attendees? & pass & says the data can't answer \\
U3 & unanswerable & How much ticket revenue did the festival make? & pass & says the data can't answer \\
U4 & unanswerable & Why exactly did people cancel? & pass & says the data can't answer \\
M1 & ambiguous & How is the partners campaign doing? & \textbf{fail} & did not ask which one \\
Q1 & data quality & Duplicate check-ins removed and denied check-ins excluded? & pass & 20 and 50 \\
Q2 & data quality & How many internal, staff or test people were excluded? & pass & 60 \\
X1 & after change & Community Groups' 14-day CPR (after logging \rs5,000)? & pass & \rs103.1 \\
\bottomrule
\end{longtable}
}

\section{REST API}
All analytics endpoints take \code{?asOf=YYYY-MM-DD}.

\begin{tabularx}{\linewidth}{@{}p{7.4cm}L@{}}
\toprule
\textbf{Endpoint} & \textbf{Purpose} \\
\midrule
\code{GET/POST /api/campaigns}, \code{GET/PATCH/DELETE /api/campaigns/:id} & Campaigns \\
\code{POST /api/campaigns/:id/metrics} & Daily metrics: JSON (one row) or CSV (many) \\
\code{GET/POST /api/tasks}, \code{PATCH/DELETE /api/tasks/:id} & Tasks \\
\code{GET /api/overview}, \code{/api/funnel}, \code{/api/insights}; \code{POST /api/scenario} & Analytics \\
\code{GET /api/activity} & Change log \\
\code{POST /api/chat} & Copilot (streamed, with trace) \\
\code{POST /api/assistant/token}, \code{POST /api/assistant/tool} & Voice: ephemeral token; server-side tools \\
\code{POST /api/admin/seed}, \code{POST /api/admin/reset} & Seed and reset (bearer token) \\
\bottomrule
\end{tabularx}

\section{More screens}
\shot{kpi-reg.png}{Registrations drill-down (REG): chart with goal line and base/stretch forecast cone, readout, and by-campaign breakdown.}{fig:kpi}
\shot{funnel.png}{Funnel: both layers, the bottleneck and the conversion matrix by campaign.}{fig:funnel}
\shot{plan.png}{Goal plan: the bottom-up model with plan vs data.}{fig:plan}
\shot{decisions.png}{Decisions: each call with its evidence and the outcome 14 days later, including the ones that got worse.}{fig:decisions}
\shot{tour-terminal.png}{First-visit tour, step 7 of 9: the copilot terminal gets its own highlighted step.}{fig:tour}

\bibliographystyle{unsrtnat}
\bibliography{refs}

\end{document}
